\section{Explicit gauges and exact arithmetic}%
\label{sec:asymex_explicit_gauge}

The compression datum of Theorem~\ref{thm:asym_multiblock} carries its
change of bond coordinates as an invertible map onto a graded coordinate
space, one summand per slot. For each worked example below, that map is an explicit
invertible matrix together with an explicit labelling of the bond
coordinates by slots and positions inside slots.

\begin{definition}[Explicit gauge and entrywise integer coercion]
    \label{def:asymex_explicit_gauge}
    \lean{MPSTensor.complexOfInt, MPSTensor.mulIntTensor, MPSTensor.gaugeOfMatrix}
    \leanok
    \uses{def:mps_tensor, def:mpdo_operator_product}
    Let $\tau$ be a bijection between the graded coordinate space of a slot
    family and the bond index set $\{0,\ldots,D_B-1\}$, and let $G$ and
    $G^{-1}$ be mutually inverse matrices in $\MN{D_B}$. The \emph{explicit
        gauge} attached to $(\tau,G)$ is the invertible map sending a bond
    vector $v$ to the graded vector $b\mapsto(Gv)(\tau b)$. For an integer
    matrix $X$, its \emph{entrywise coercion} is the complex matrix with the
    same entries; for integer matrix product operator tensors $M$ and $N$,
    their integer bond-space product is
    $(M\mathbin{\cdot}N)^{ik}=\sum_jM^{ij}\otimes N^{jk}$.
\end{definition}

\begin{theorem}[Conjugation, coercion and the compression pairs of a slot]
    \label{thm:asymex_explicit_gauge}
    % Principal declaration: MPSTensor.conjMatrix_gaugeOfMatrix.
    \lean{MPSTensor.conjMatrix_gaugeOfMatrix, MPSTensor.complexOfInt_mul,
        MPSTensor.mulTensor_complexOfRing, MPSTensor.mulTensor_smul_complexOfRing,
        MPSTensor.MultiBlockCompression.left_gaugeOfMatrix,
        MPSTensor.MultiBlockCompression.right_gaugeOfMatrix}
    \leanok
    \uses{def:asymex_explicit_gauge, thm:asym_multiblock}
    Conjugation by the explicit gauge attached to $(\tau,G)$ sends a bond
    matrix $A$ to the matrix $GAG^{-1}$ relabelled along $\tau$. The
    entrywise coercion is multiplicative, and the bond-space product of two
    integer tensors each rescaled by one scalar $c$ is $c^2$ times the
    entrywise coercion of their integer bond-space product. If a compression
    datum uses the explicit gauge attached to $(\tau,G)$, then the
    compression map $W_s$ out of the bond space consists of the rows of $G$
    at the coordinates of the slot $s$, and the compression map $V_s$ into
    the bond space consists of the columns of $G^{-1}$ at those same
    coordinates.
\end{theorem}

\begin{proof}
    \leanok
    The matrix of the explicit gauge is $G$ relabelled along $\tau$ in its
    rows, and the matrix of its inverse is $G^{-1}$ relabelled along $\tau$
    in its columns, so conjugation is the stated matrix conjugation followed
    by the relabelling. Multiplicativity of the coercion is the fact that a
    sum of products of integers is the same computed in $\Z$ or in $\C$, and
    the statement about the bond-space product follows by expanding both
    sides entrywise. The two compression maps are the composites of the
    coordinate embedding and projection of the slot with the gauge and its
    inverse; evaluating the matrix products at a single graded coordinate
    leaves one row of $G$ and one column of $G^{-1}$.
\end{proof}

\begin{definition}[One-slot compression data]%
    \label{def:asymex_one_slot}
    \lean{MPSTensor.stackedInt, MPSTensor.oneSlot, MPSTensor.oneSlotMem,
        MPSTensor.MultiBlockCompression.ofConjMatrix,
        MPSTensor.MultiBlockCompression.ofConjInt,
        MPSTensor.MultiBlockCompression.ofEq,
        MPSTensor.MultiBlockCompression.ofScalarFlagFour}
    \leanok
    \uses{def:asymex_explicit_gauge, thm:asym_multiblock}
    For a stacked product of two integer matrix product operator tensors,
    read over the pair alphabet, the letter at $a=do+i$ is
    $(M\mathbin{\cdot}N)^{oi}$. A \emph{one-slot compression datum} for a
    tensor $B$ onto a single target $C$ is the data of
    Theorem~\ref{thm:asym_multiblock} for the one-element slot set: an
    explicit gauge $(\tau,G)$ conjugating $B$ into a matrix that is block
    upper triangular for a flag with $C$ at one diagonal position. Two
    shapes recur in the examples below: a conjugation $GB^iG^{-1}=C^i$ for
    an invertible $G$ gives such a datum with no zero slots ($z=0$); and,
    for a four-dimensional $B$ and a one-dimensional $C$, an integer gauge
    sending every letter to a matrix with $C$ at the third diagonal
    position of a four-step flag and zero at the other three gives such a
    datum with three zero slots ($z=3$).
\end{definition}

\begin{theorem}[The remainder of a conjugate datum vanishes; word traces
        of a one-slot datum]%
    \label{thm:asymex_one_slot}
    \lean{MPSTensor.MultiBlockCompression.remainder_ofConjMatrix,
        MPSTensor.MultiBlockCompression.remainder_ofEq,
        MPSTensor.MultiBlockCompression.remainder_ofConjInt,
        MPSTensor.MultiBlockCompression.trace_evalWord_oneSlot,
        MPSTensor.MultiBlockCompression.trace_evalWord_oneSlot_smul}
    \leanok
    \uses{def:asymex_one_slot}
    The remainder of a conjugate one-slot datum vanishes identically: the
    compression is an exact change of coordinates. For every one-slot
    datum, the trace of every nonempty word of $B$ equals the trace of the
    corresponding word of $C$, or, if the target is carried with a scalar
    weight $c$, that power of $c$ times the trace of the word of the
    unweighted target.
\end{theorem}

\begin{proof}
    \leanok
    \uses{thm:asymex_explicit_gauge, prop:asym_compression_trace}
    Immediate from Theorem~\ref{thm:asymex_explicit_gauge} and
    Proposition~\ref{prop:asym_compression_trace}, specialized to a single
    slot.
\end{proof}

\begin{theorem}[Integer certificates for normality and for vanishing
        sitewise intertwiners]%
    \label{thm:asymex_int_certificates}
    \lean{MPSTensor.isNBlkInjective_two_of_int,
        MPSTensor.right_intertwiner_eq_zero_of_int,
        MPSTensor.left_intertwiner_eq_zero_of_int}
    \leanok
    \uses{def:normal}
    If every matrix unit is, up to a common nonzero integer factor, an
    integer combination of length-two words of an integer tensor, the
    length-two words span the full matrix algebra. If the rows,
    respectively columns, of the letters occurring in a sitewise
    intertwiner system form an integer matrix with an integer left inverse
    up to a nonzero factor, the only solution of that system is zero.
\end{theorem}

\begin{proof}
    \leanok
    Clearing the nonzero integer factor by scalar division after coercion
    to $\C$ reduces each statement to the corresponding identity between
    integer matrices, a finite verification.
\end{proof}

\section{Hidden symmetry-broken sectors}%
\label{sec:asymex_ghz}

The following example describes two broken sectors of a $\Z_2$ symmetry
by a periodic vector family generated by the sum of two inequivalent
one-dimensional normal blocks, one per sector. A tensor generating that same family need
not display the sectors: they may occupy non-adjacent positions, in an order
different from the target's, coupled through a zero block. This is the
smallest example with more than one target slot, and it serves as a warm-up
for the fusion examples below.

\begin{definition}[The two sectors and the four-dimensional source]
    \label{def:asymex_ghz}
    \lean{MPSTensor.ghzC, MPSTensor.ghzB}
    \leanok
    \uses{def:mps_tensor}
    The target is the pair of one-dimensional tensors
    $C_0=\bigl((1),(0)\bigr)$ and $C_1=\bigl((0),(1)\bigr)$ on the alphabet
    $\{0,1\}$, whose direct sum is the Greenberger--Horne--Zeilinger tensor
    $\diag(1,0)$, $\diag(0,1)$. The source is the bond-dimension-four
    tensor
    \begin{align}
        B^0
         & =\begin{pmatrix}
                0 & 1 & 2 & 0 \\
                0 & 0 & 1 & 1 \\
                0 & 0 & 1 & 3 \\
                0 & 0 & 0 & 0
            \end{pmatrix},
         &
        B^1
         & =\begin{pmatrix}
                1 & 1 & 0 & 2 \\
                0 & 0 & 2 & 0 \\
                0 & 0 & 0 & 1 \\
                0 & 0 & 0 & 0
            \end{pmatrix}.
        \notag
    \end{align}
\end{definition}

\begin{center}
    % Formula: the periodic trace tr(B^w) of the bond-dimension-four
    % source of this section.
    % Source: Definition~\ref{def:asymex_ghz}.
    % Ink-to-index map: each node is the source tensor $B$; the upper
    % legs carry the letters of the word.
    % Contraction set: the horizontal virtual legs, closed around the
    % ring by the trace; no leg is left open.
    % Boundary signature: (virtual-west=0 | virtual-east=0 |
    % physical-up=3 | physical-down=0) for the finite schematic shown.
    \tenkzkernel
    \begin{tenkz}[cols=4, boundary=periodic, physical=up]
        \tn[label pos=135, ports={90:physical:$i_1$}]{B} &
        \tn[label pos=135]{B} & \tn[skin=dots]{} &
        \tn[label pos=135, ports={90:physical:$i_N$}]{B}
    \end{tenkz}
\end{center}

\begin{theorem}[Compression of the two sectors]
    \label{thm:asymex_ghz_compression}
    % Principal declaration: MPSTensor.ghzSectors_isReduction.
    \lean{MPSTensor.ghzSectors_compression, MPSTensor.ghzSectors_isReduction,
        MPSTensor.ghzSectors_dim_eq}
    \leanok
    \uses{def:asymex_ghz, thm:asym_multiblock}
    The source of Definition~\ref{def:asymex_ghz} carries the block
    triangular form of Theorem~\ref{thm:asym_multiblock} for the slot family
    $(C_0,C_1)$ with $z=2$ zero slots, the four diagonal positions carrying,
    in order, the sector $C_1$, a zero slot, the sector $C_0$ and a second
    zero slot. The compression pair of each sector satisfies $W_sV_s=\id$
    and $W_sB^wV_s=C_s^w$ for every word, and $4=1+1+2$.
\end{theorem}

\begin{proof}
    \leanok
    \uses{thm:asymex_explicit_gauge}
    The gauge is the relabelling of the four bond coordinates by the four
    slots, with no change of basis. The three matrix conditions of the
    compression datum, block triangularity and the identification of the
    matched and unmatched diagonal blocks, are then a finite verification
    over the integer entries of $B^0$ and $B^1$. The compression pairs and
    the dimension count are clauses (iv), (v) and (vii) of
    Theorem~\ref{thm:asym_multiblock}.
\end{proof}

\begin{theorem}[Periodic data of the four-dimensional source]
    \label{thm:asymex_ghz_trace}
    \lean{MPSTensor.ghzB_trace_evalWord_eq_sum}
    \leanok
    \uses{def:asymex_ghz}
    For every nonempty word $w$, $\tr(B^w)=\tr(C_0^w)+\tr(C_1^w)$.
\end{theorem}

\begin{proof}
    \leanok
    \uses{thm:asymex_ghz_compression, prop:asym_compression_trace}
    This is Proposition~\ref{prop:asym_compression_trace} applied to the
    compression datum of Theorem~\ref{thm:asymex_ghz_compression}, whose
    slot set has two members.
\end{proof}

\begin{theorem}[The first sector has no sitewise intertwiner]
    \label{thm:asymex_ghz_no_intertwiner}
    \lean{MPSTensor.ghzC0_right_intertwiner_eq_zero,
        MPSTensor.ghzC0_left_intertwiner_eq_zero}
    \leanok
    \uses{def:asymex_ghz}
    If $v\in\C^4$ satisfies $B^0v=v$ and $B^1v=0$, then $v=0$; if
    $u\in\C^4$ satisfies $uB^0=u$ and $uB^1=0$, then $u=0$. Both sitewise
    intertwiner spaces for the sector $C_0$ therefore vanish, and the
    word-level compression of Theorem~\ref{thm:asymex_ghz_compression} is
    the strongest local relation available for that sector.
\end{theorem}

\begin{proof}
    \leanok
    The last row of $B^0$ vanishes, so the last coordinate of $B^0v=v$
    gives $v_3=0$; the second coordinate of $B^1v=0$ gives $v_2=0$; and the
    second and first coordinates of $B^0v=v$ then force $v_1=0$ and
    $v_0=0$. For the covector, the first and third coordinates of $uB^1=0$
    give $u_0=0$ and $u_1=0$, its last coordinate gives $u_2=0$, and the
    last coordinate of $uB^0=u$ gives $u_3=0$.
\end{proof}

\section{A repeated block with a nontrivial extension}%
\label{sec:asymex_repeated}

Periodic traces determine the multiplicity with which a simple block occurs,
but not the extension carrying the copies. The following warm-up shows that
the conclusion of the symmetric fundamental theorem, an invertible gauge
relating source and target, is false for such a source, not merely
inapplicable.

\begin{definition}[The scalar block and its Jordan extension]
    \label{def:asymex_repeated}
    \lean{MPSTensor.repA, MPSTensor.repC, MPSTensor.repB}
    \leanok
    \uses{def:mps_tensor}
    The target is the scalar normal tensor $A^0=(1)$, $A^1=(2)$, taken
    twice, and the source is
    \begin{align}
        B^0
         & =\begin{pmatrix}1&1\\0&1\end{pmatrix},
         &
        B^1
         & =\begin{pmatrix}2&0\\0&2\end{pmatrix}.
        \notag
    \end{align}
\end{definition}

\begin{center}
    % Formula: the periodic trace tr(B^w) of the bond-dimension-two
    % Jordan source of this section.
    % Source: Definition~\ref{def:asymex_repeated}.
    % Ink-to-index map: each node is the source tensor $B$; the upper
    % legs carry the letters of the word.
    % Contraction set: the horizontal virtual legs, closed around the
    % ring by the trace; no leg is left open.
    % Boundary signature: (virtual-west=0 | virtual-east=0 |
    % physical-up=3 | physical-down=0) for the finite schematic shown.
    \tenkzkernel
    \begin{tenkz}[cols=4, boundary=periodic, physical=up]
        \tn[label pos=135, ports={90:physical:$i_1$}]{B} &
        \tn[label pos=135]{B} & \tn[skin=dots]{} &
        \tn[label pos=135, ports={90:physical:$i_N$}]{B}
    \end{tenkz}
\end{center}

\begin{theorem}[Compression onto the two copies]
    \label{thm:asymex_repeated_compression}
    % Principal declaration: MPSTensor.repeatedBlock_isReduction.
    \lean{MPSTensor.repeatedBlock_compression, MPSTensor.repeatedBlock_isReduction,
        MPSTensor.repB_trace_evalWord_eq_two_mul}
    \leanok
    \uses{def:asymex_repeated, thm:asym_multiblock}
    The source of Definition~\ref{def:asymex_repeated} carries the block
    triangular form of Theorem~\ref{thm:asym_multiblock} for the two copies
    of $A$ with no zero slots, each coordinate pair compresses its copy,
    $W_sV_s=\id$ and $W_sB^wV_s=A^w$ for every word, and
    $\tr(B^w)=2\tr(A^w)$ for every nonempty word.
\end{theorem}

\begin{proof}
    \leanok
    \uses{prop:asym_compression_trace}
    Both matrices are upper triangular with both diagonal entries equal to
    the corresponding scalar of $A$, which is the whole content of the three
    matrix conditions; the compression pairs are clauses (iv) and (v) of
    Theorem~\ref{thm:asym_multiblock}, and the trace identity is
    Proposition~\ref{prop:asym_compression_trace} with two equal slots.
\end{proof}

\begin{theorem}[No invertible gauge onto the repeated target]
    \label{thm:asymex_repeated_not_gauge}
    \lean{MPSTensor.not_gaugeEquiv_directSum_repA}
    \leanok
    \uses{def:asymex_repeated, def:gauge_equiv}
    There is no invertible $X\in\GL_2(\C)$ with
    $B^i=X\diag(A^i,A^i)X^{-1}$ for both letters $i$.
\end{theorem}

\begin{proof}
    \leanok
    For $i=0$ the matrix $\diag(A^0,A^0)$ is the identity, so any such
    conjugation would give $B^0=\id$, contradicting the nonzero
    off-diagonal entry of $B^0$.
\end{proof}

\section{The square of the CZX $\Z_2$ symmetry}%
\label{sec:asymex_czx}

On a periodic chain of $L$ qubits the symmetry of CZX type is
\begin{align}
    U_L & =\Bigl(\prod_{i=1}^{L}CZ_{i,i+1}\Bigr)X^{\otimes L},
    \label{eq:asymex_czx_operator}
\end{align}
a matrix product operator of bond dimension two with tensor
$M^{ij}=\delta_{i,1\oplus j}T_j$. This is the CZX example of an anomalous
$\Z_2$ symmetry. Its square is $U_L^2=(-1)^L\id^{\otimes L}$, so the stacked product of its
tensor with itself has, at every positive length, the periodic data of the
one-dimensional identity tensor carried with the weight $-1$. The fusion is
therefore beyond the unit-weight setting. The sitewise theory of
\cite[Appendix~A]{GarreRubio2023MPOSym} imposes a closedness requirement
for arbitrary boundary conditions. The compression and intertwiner
calculations below concern this explicit tensor, not a classification of
representations by their three-cocycle.

\begin{definition}[The bond-two symmetry tensor and its square]
    \label{def:asymex_czx_tensors}
    \lean{CZXCompression.czxTensor, CZXCompression.czxMPS,
        CZXCompression.identityTensor, CZXCompression.identityMPS,
        CZXCompression.czxSquare, CZXCompression.czxSquareTarget}
    \leanok
    \uses{def:mpo_tensor, def:mpo_to_mps, def:mpdo_operator_product}
    The symmetry tensor is $M^{ij}=\delta_{i,1\oplus j}T_j$ with
    \begin{align}
        T_0 & =\begin{pmatrix}1&0\\1&0\end{pmatrix},
            &
        T_1 & =\begin{pmatrix}0&1\\0&-1\end{pmatrix};
        \notag
    \end{align}
    the identity tensor is $\delta^{ij}$ of bond dimension one. Read over
    the alphabet of index pairs, the symmetry tensor has the four matrices
    $0$, $T_1$, $T_0$, $0$ in the order $(0,0)$, $(0,1)$, $(1,0)$, $(1,1)$.
    The source is the stacked product
    $B^{ij}=\sum_mM^{im}\otimes M^{mj}$ of bond dimension four, and the
    single target is $-\delta$.
\end{definition}

\begin{center}
    % Formula: B^{ij}=\sum_mM^{im}\otimes M^{mj}, the stacked product of
    % Definition~\ref{def:asymex_czx_tensors}.
    % Source: Definition~\ref{def:asymex_czx_tensors}.
    % Ink-to-index map: both rows are the bond-two symmetry tensor $M$;
    % the shared vertical pairleg at each site is the summed index $m$.
    % Contraction set: the vertical pairlegs (index $m$) at every site
    % and the horizontal virtual bonds within each row; the physical
    % legs $(i,j)$ remain open.
    % Boundary signature: (virtual-west=0 | virtual-east=0 |
    % physical-up=3 | physical-down=3) on both sides.
    \tenkzkernel
    \begin{tenkz}[rows={op,op}, cols=4, boundary=periodic]
        \tn[skin=mpo, ports={90:physical:$i_1$, 270:physical}]{M} &
        \tn[skin=mpo, ports={90:physical:$i_2$, 270:physical}]{M} &
        \tn[skin=dots, ports={180:virtual, 0:virtual}]{} &
        \tn[skin=mpo, ports={90:physical:$i_N$, 270:physical}]{M} \\
        \tn[skin=mpo, ports={90:physical, 270:physical:$j_1$}]{M} &
        \tn[skin=mpo, ports={90:physical, 270:physical:$j_2$}]{M} &
        \tn[skin=dots, ports={180:virtual, 0:virtual}]{} &
        \tn[skin=mpo, ports={90:physical, 270:physical:$j_N$}]{M}
    \end{tenkz}
    \quad $=$ \quad
    \begin{tenkz}[rows={op}, cols=4, boundary=periodic]
        \tn[skin=mpo, ports={90:physical:$i_1$, 270:physical:$j_1$}]{B} &
        \tn[skin=mpo, ports={90:physical:$i_2$, 270:physical:$j_2$}]{B} &
        \tn[skin=dots, ports={180:virtual, 0:virtual}]{} &
        \tn[skin=mpo, ports={90:physical:$i_N$, 270:physical:$j_N$}]{B}
    \end{tenkz}
\end{center}

\begin{theorem}[Normality of the symmetry tensor and of the target]
    \label{thm:asymex_czx_normal}
    \lean{CZXCompression.czxMPS_isNormal, CZXCompression.czxSquareTarget_isNormal}
    \leanok
    \uses{def:asymex_czx_tensors, def:normal}
    The bond-two symmetry tensor is normal: its length-two words span
    $\MN{2}$. The target $-\delta$ is normal: its bond dimension is one and
    one of its matrices is the nonzero scalar $-1$.
\end{theorem}

\begin{proof}
    \leanok
    \uses{thm:asymex_normal_units}
    Each of the four matrix units of $\MN{2}$ is a half-sum or a
    half-difference of two of the four products $T_aT_b$, which is a finite
    verification over integer matrices; the span of the length-two words is
    therefore all of $\MN{2}$. For the target, its matrix at the letter
    $(0,0)$ is the nonzero scalar $-1$, so
    Theorem~\ref{thm:asymex_normal_units} applies with the single position
    of the one-by-one matrix algebra.
\end{proof}

\begin{theorem}[Compression of the squared symmetry]
    \label{thm:asymex_czx_compression}
    % Principal declaration: CZXCompression.czxSquare_isReduction.
    \lean{CZXCompression.czxSquare_compression, CZXCompression.czxSquareLeft,
        CZXCompression.czxSquareRight, CZXCompression.czxSquare_isReduction,
        CZXCompression.czxSquareLeft_mul_right, CZXCompression.czxSquare_dim_eq,
        CZXCompression.czxSquare_evalWord_remainder_eq_zero}
    \leanok
    \uses{def:asymex_czx_tensors, thm:asym_multiblock}
    The stacked product of Definition~\ref{def:asymex_czx_tensors} carries
    the block triangular form of Theorem~\ref{thm:asym_multiblock} for the
    single slot $-\delta$ with $z=3$ zero slots and $4=1+3$. The
    compression pair is
    \begin{align}
        W & =(0,0,-1,0),
          &
        V & =(1,1,-1,-1)^{\mathsf T},
        \notag
    \end{align}
    with $WV=1$ and $WB^wV=(-\delta)^w$ for every word, and the remainder
    $R^{ij}=B^{ij}-V(-\delta)^{ij}W$ satisfies
    $R^{i_1j_1}R^{i_2j_2}R^{i_3j_3}R^{i_4j_4}=0$.
\end{theorem}

\begin{proof}
    \leanok
    \uses{thm:asymex_explicit_gauge}
    The change of bond coordinates has as its columns, in the order of the
    flag, the line $(1,0,0,-1)$ annihilated by every letter, the vector
    $(1,1,-1,-1)$ spanning the image of every letter modulo that line, and
    two arbitrary completions. In these coordinates every letter is block
    upper triangular with the entry $-1$ in the second diagonal position and
    zero in the other three, which is a finite verification over integer
    matrices. The explicit pair is
    Theorem~\ref{thm:asymex_explicit_gauge} applied to the second
    coordinate, and the vanishing of a product of four remainder matrices is
    clause (vi) of Theorem~\ref{thm:asym_multiblock} with $|S|+z=4$.
\end{proof}

\begin{theorem}[Periodic data of the squared symmetry]
    \label{thm:asymex_czx_trace}
    \lean{CZXCompression.czxSquare_trace_evalWord}
    \leanok
    \uses{def:asymex_czx_tensors}
    For every nonempty word $w$ over the alphabet of index pairs,
    $\tr(B^w)=(-1)^{|w|}\tr(\delta^w)$: this is
    $U_L^2=(-1)^L\id^{\otimes L}$ at the level of periodic coefficients.
\end{theorem}

\begin{proof}
    \leanok
    \uses{thm:asymex_czx_compression, prop:asym_compression_trace}
    Proposition~\ref{prop:asym_compression_trace} applied to the datum of
    Theorem~\ref{thm:asymex_czx_compression} gives
    $\tr(B^w)=\tr\bigl((-\delta)^w\bigr)$, and pulling the scalar out of the
    word evaluation contributes the factor $(-1)^{|w|}$.
\end{proof}

\begin{theorem}[Both sitewise intertwiner spaces vanish]
    \label{thm:asymex_czx_no_intertwiner}
    \lean{CZXCompression.czxSquare_right_intertwiner_eq_zero,
        CZXCompression.czxSquare_left_intertwiner_eq_zero}
    \leanok
    \uses{def:asymex_czx_tensors}
    If $v\in\C^4$ satisfies $B^{ij}v=(-\delta)^{ij}v$ for every index pair,
    then $v=0$; if $u\in\C^4$ satisfies $uB^{ij}=(-\delta)^{ij}u$ for every
    index pair, then $u=0$.
\end{theorem}

\begin{proof}
    \leanok
    Only the letters $(0,0)$ and $(1,1)$ of the stacked tensor are nonzero,
    and each is a rank-one matrix with an explicit integer kernel and image.
    Comparing the four coordinates of the two equations for those two
    letters gives, in turn, $v_1=v_2=v_0=v_3=0$, and the transposed
    comparison gives $u_0=u_1=u_3=u_2=0$.
\end{proof}

\begin{remark}[Failure of closure under conjugate transposition]
    \label{rem:asymex_czx_anomaly}
    \uses{cor:asym_anomaly_not_star}
    Theorem~\ref{thm:asymex_czx_no_intertwiner} is what
    Corollary~\ref{cor:asym_anomaly_not_star} converts into the statement
    that the algebra generated by the matrices of the stacked tensor is not
    closed under conjugate transposition, although the symmetry operator
    itself is unitary. The obstruction to a sitewise fusion pair is
    therefore algebraic and not a defect of the particular coordinates
    chosen here. This algebraic implication does not identify nonsplitting
    with an anomaly class. Determining that class requires a separate
    computation of the fusion-tensor associator and its cohomology class.
\end{remark}

\section{A non-simple symmetry family with an oscillating coefficient}%
\label{sec:asymex_czx_plus}

Adding the identity to the symmetry of Section~\ref{sec:asymex_czx} gives
the non-simple family $O_L=\id^{\otimes L}+U_L$, generated by the
bond-dimension-three tensor $N=\delta\oplus M$ whose two blocks are normal
and inequivalent. Its square is
\begin{align}
    O_L^2 & =\bigl(1+(-1)^L\bigr)\id^{\otimes L}+2\,U_L,
    \label{eq:asymex_czx_plus_square}
\end{align}
an operator family whose identity structure coefficient oscillates with the length
while remaining an integer. The compression resolves the oscillation
exactly as Theorem~\ref{thm:asym_power_sum_coefficients} predicts: the
identity block occurs twice, with the weights $+1$ and $-1$, and
$1+(-1)^L$ is the length-$L$ power sum of the weight multiset
$\{+1,-1\}$. This is an example of the general power-sum mechanism,
not of the positive diagonal matrices in the CPSV characterization
of MPDO renormalization fixed points.

\begin{definition}[The non-simple family and its square]
    \label{def:asymex_czx_plus}
    \lean{CZXCompression.plusTensor, CZXCompression.czxPlusIdentity,
        CZXCompression.plusBlockDim, CZXCompression.plusTargets}
    \leanok
    \uses{def:asymex_czx_tensors, def:mpdo_operator_product}
    The generating tensor is $N^{ij}=\delta^{ij}\oplus M^{ij}$ of bond
    dimension three, the source is the stacked product
    $B^{ij}=\sum_mN^{im}\otimes N^{mj}$ of bond dimension nine, and the four
    target slots are $\delta$, $-\delta$, $M$ and $M$, of bond dimensions
    $1$, $1$, $2$ and $2$.
\end{definition}

\begin{center}
    % Formula: B^{ij}=\sum_mN^{im}\otimes N^{mj}, the stacked product of
    % Definition~\ref{def:asymex_czx_plus}.
    % Source: Definition~\ref{def:asymex_czx_plus}.
    % Ink-to-index map: both rows are the bond-three generating tensor
    % $N=\delta\oplus M$; the shared vertical pairleg at each site is
    % the summed index $m$.
    % Contraction set: the vertical pairlegs (index $m$) at every site
    % and the horizontal virtual bonds within each row; the physical
    % legs $(i,j)$ remain open.
    % Boundary signature: (virtual-west=0 | virtual-east=0 |
    % physical-up=3 | physical-down=3) on both sides.
    \tenkzkernel
    \begin{tenkz}[rows={op,op}, cols=4, boundary=periodic]
        \tn[skin=mpo, ports={90:physical:$i_1$, 270:physical}]{N} &
        \tn[skin=mpo, ports={90:physical:$i_2$, 270:physical}]{N} &
        \tn[skin=dots, ports={180:virtual, 0:virtual}]{} &
        \tn[skin=mpo, ports={90:physical:$i_N$, 270:physical}]{N} \\
        \tn[skin=mpo, ports={90:physical, 270:physical:$j_1$}]{N} &
        \tn[skin=mpo, ports={90:physical, 270:physical:$j_2$}]{N} &
        \tn[skin=dots, ports={180:virtual, 0:virtual}]{} &
        \tn[skin=mpo, ports={90:physical, 270:physical:$j_N$}]{N}
    \end{tenkz}
    \quad $=$ \quad
    \begin{tenkz}[rows={op}, cols=4, boundary=periodic]
        \tn[skin=mpo, ports={90:physical:$i_1$, 270:physical:$j_1$}]{B} &
        \tn[skin=mpo, ports={90:physical:$i_2$, 270:physical:$j_2$}]{B} &
        \tn[skin=dots, ports={180:virtual, 0:virtual}]{} &
        \tn[skin=mpo, ports={90:physical:$i_N$, 270:physical:$j_N$}]{B}
    \end{tenkz}
\end{center}

\begin{theorem}[Compression of the squared non-simple family]
    \label{thm:asymex_czx_plus_compression}
    % Principal declaration: CZXCompression.czxPlusIdentity_isReduction.
    \lean{CZXCompression.czxPlusIdentity_compression,
        CZXCompression.czxPlusIdentity_isReduction,
        CZXCompression.czxPlusIdentity_left_mul_right_of_ne,
        CZXCompression.czxPlusIdentity_dim_eq,
        CZXCompression.czxPlusIdentity_evalWord_remainder_eq_zero}
    \leanok
    \uses{def:asymex_czx_plus, thm:asym_multiblock}
    The source of Definition~\ref{def:asymex_czx_plus} carries the block
    triangular form of Theorem~\ref{thm:asym_multiblock} for the four
    slots with $z=3$ zero slots and $9=1+1+2+2+3$. The compression pairs
    are mutually biorthogonal, $W_sV_t=\id$ for $s=t$ and $0$ otherwise,
    they satisfy $W_sB^wV_s=C_s^w$ for every word, and a product of seven
    remainder matrices vanishes.
\end{theorem}

\begin{proof}
    \leanok
    \uses{thm:asymex_explicit_gauge}
    Every letter of $N$ is block diagonal for the splitting
    $\C^3=\C\oplus\C^2$, so every letter of the stacked tensor is block
    diagonal for the induced four-fold splitting of $\C^3\otimes\C^3$.
    Three of those four summands already carry the simple blocks $\delta$,
    $M$ and $M$; the fourth is the stacked tensor of
    Section~\ref{sec:asymex_czx}, whose flag contributes the weight $-1$
    together with the three zero slots. Assembling the corresponding change
    of bond coordinates and checking block triangularity together with the
    matched and unmatched diagonal blocks is a finite verification over
    integer matrices. The dimension count and the nilpotency order are
    clauses (vii) and (vi) of Theorem~\ref{thm:asym_multiblock} with
    $|S|+z=7$.
\end{proof}

\begin{theorem}[The oscillating structure coefficient]
    \label{thm:asymex_czx_plus_trace}
    \lean{CZXCompression.czxPlusIdentity_trace_evalWord}
    \leanok
    \uses{def:asymex_czx_plus}
    For every nonempty word $w$ over the alphabet of index pairs,
    \begin{align}
        \tr(B^w) & =\bigl(1+(-1)^{|w|}\bigr)\tr(\delta^w)+2\tr(M^w).
        \label{eq:asymex_czx_plus_trace}
    \end{align}
\end{theorem}

\begin{proof}
    \leanok
    \uses{thm:asymex_czx_plus_compression, prop:asym_compression_trace}
    Proposition~\ref{prop:asym_compression_trace} applied to the datum of
    Theorem~\ref{thm:asymex_czx_plus_compression} writes $\tr(B^w)$ as the
    sum of the four slot traces; pulling the weight $-1$ out of the second
    slot contributes $(-1)^{|w|}\tr(\delta^w)$ and the two equal slots
    carrying $M$ contribute $2\tr(M^w)$.
\end{proof}

